\section{Conclusions}
% C1: coupled formulation and physical evaluation
The coupled formulations connect five-reaction, nine-species liquid equilibrium and ePC-SAFT bulk-phase properties to nonisothermal MEA absorber balances.
ePC-SAFT supplies the density, enthalpies, heat capacities, phase and reaction equilibrium, \COtwo{} solubility and water fugacities these balances use, while surface tension, transport, kinetics and packing performance remain correlations.
No thermodynamic parameter is fitted to the seven NCCC one-bed cases, so the capture and temperature comparisons are predictions.

% C2: twelve-state film comparison
In the twelve-state conserved formulation, a finite-rate reactive film with the ePC-SAFT Henry coefficient predicts capture with a mean absolute error of \num{3.55} percentage points at 129 axial nodes, against \num{4.52} for an equilibrium-manifold film and \num{7.51} for conventional enhancement.
The \NtwoO{}-analogy solubility gives \num{3.61} points, so these cases do not distinguish the two solubility bases.
Every film overpredicts Case 6C, Case 7C capture is not resolved with respect to the axial mesh, and the calculated liquid temperatures exceed the Case 7C packing temperatures by up to \qty{24}{K}.

% C3: nine-species numerical behavior and property sources
The nine-species BVP formulation with conventional enhancement overpredicts all seven cases, with a mean absolute error of \num{5.86} percentage points.
Its refined Case 3C calculation predicts \qty{92.91}{\percent} capture against \qty{89.50}{\percent} observed; mesh refinement changes capture by \num{0.005} points, and one coarse solution requires about \qty{17}{s} of BVP wall time.
Replacing empirical density and calorics with ePC-SAFT lowers capture in every case and raises the packing-temperature error; in Case 3C the capture change comes mostly from the lower ePC-SAFT water vaporization enthalpy.

% C4: thermodynamic, transport and operating insights
The MEA--water interaction produces the largest capture response among the tested thermodynamic inputs and the liquid-side transfer coefficient the largest among the transport inputs, but neither closes the Case 3C capture gap within the tested ranges.
Lean loadings of 0.225 and 0.275 mol \COtwo{} per mol MEA change capture by \num{+2.41} and \num{-3.92} points relative to 0.25, identifying solvent loading as the most influential operating variable under the tested conditions.

The temperature comparison is qualitative because the source does not identify the measured phase; the conventional enhancement and the transport correlations used beyond their source ranges constrain predictive interpretation.

% C5: contribution and forward program
The formulations give a direct connection from component-level molecular and ionic descriptions to column-scale observables and operating response.
The minimum Case 3C solvent flow for \qty{90}{\percent} capture from the twelve-state model moves toward the reference flow with mesh refinement, leaving its mesh-independent value unresolved.
A direct ePC-SAFT/eNRTL comparison and tests with other amines would evaluate how thermodynamic parameter reuse contributes to solvent selection and absorber design.

\clearpage
\nolinenumbers
\raggedbottom
